\section*{Hoofdstuk \ref{ch:b3:milky-way-galaxies} --- De Melkweg en sterrenstelsels}

\begin{solution}{exo:b3:milky-way-galaxies:1}
(a)~$2\pi R_0 = 2\pi\times\qty{2.523e17}{km} = \qty{1.585e18}{km}$; bij
\qty{229}{km.s^{-1}} duurt dit \qty{6.92e15}{s}, dat is
\qty{219}{Myr}. (b)~$4540/219 = 20.7$: ongeveer twintig omwentelingen.
\end{solution}

\begin{solution}{exo:b3:milky-way-galaxies:2}
(a)~$\qty{1}{au}/\qty{1}{yr} = \num{1.496e8}/\num{3.156e7} =
\qty{4.74}{km.s^{-1}}$. (b)~$v = 4.74\times6.411\times8.178 =
\qty{248.5}{km.s^{-1}}$, binnen de waarde van het overzicht
\qty[separate-uncertainty=true]{248 \pm 3}{km.s^{-1}}. (c)~$248.5 - 11 = \qty{237.5}{km.s^{-1}}$,
ongeveer \qty{8}{km.s^{-1}} boven de \qty{229}{km.s^{-1}} van de rode reuzen:
de twee berusten op verschillende waarden van $R_0$ (8.178 en \qty{8.122}{kpc}),
en de LSR zelf kan tot ongeveer \qty{15}{km.s^{-1}} meebewegen ten opzichte van
een werkelijk cirkelvormige baan.
\end{solution}

\begin{solution}{exo:b3:milky-way-galaxies:3}
(a)~$n = 10(1 - 0.6) = 4$: E4. (b)~$b/a = 1 - 7/10 = 0.3$. (c)~De klasse
gebruikt de schijnbare asverhouding: een afgeplatte ellipsoïde, gezien langs zijn korte
as heeft een cirkelvormige omtrek en wordt als E0 ingedeeld.
\end{solution}

\begin{solution}{exo:b3:milky-way-galaxies:4}
$d = 10^{\mu/5 + 1}$~pc: de LMC \qty{49.6}{kpc}, Andromeda
\qty{783}{kpc}, NGC~4258 \qty{7.58}{Mpc}. Een ster van de Wolk heeft een
parallaxis van $1/\qty{49.6}{kpc} = \qty{0.020}{mas}$, de grootte van de
onzekerheid van de beste catalogus voor heldere sterren: één parallaxis geeft helemaal geen
afstand, en de Wolk wordt met andere middelen vastgesteld (bedekkingsdubbelsterren),
waaruit de cepheïdeschaal begint.
\end{solution}

\begin{solution}{exo:b3:milky-way-galaxies:5}
(a)~$\omega_0 = A - B = \qty{27.2}{km.s^{-1}.kpc^{-1}}$ en $V_0 =
27.2\times8.178 = \qty{222}{km.s^{-1}}$. (b)~$V'_0 = -(A + B) =
\qty{-3.4}{km.s^{-1}.kpc^{-1}}$: een geleidelijk dalende \omterm{def:b3:milky-way-galaxies:rotation-curve}{rotatiekromme}. (c)~Beide zijn
onder de waarden van de rode reuzen (229 en \qty{-1.7}{km.s^{-1}} per
kiloparsec); de auteurs van die kromme vonden hetzelfde, een $V_0$ ongeveer twee
standaardafwijkingen lager door $R_0$ en deze \omterm{def:b3:milky-way-galaxies:rotation-curve}{Oort-constanten}. De formules van Oort
gaan uit van cirkelvormige banen in een axiaalsymmetrische schijf; niet-nul $C$ en
$K$ tonen aan dat het lokale snelheidsveld wordt verstoord door de balk en de
spiraalarmen, en de lokale $A$ en $B$ bevatten die verstoringen.
\end{solution}

\begin{solution}{exo:b3:milky-way-galaxies:6}
(a)~$\rho = V^2/(4\pi GR^2) = 229^2/(4\pi\times\num{4.301e-6}\times
8.178^2) = \num{1.45e7}$ zonsmassa's per kubieke kiloparsec, of
\num{0.0145} per kubieke parsec. (b)~$M_\odot c^2 = \num{1.9884e30}\times
(\num{2.998e8})^2 = \qty{1.787e47}{J} = \qty{1.115e57}{GeV}$, en
$\qty{1}{pc^3} = (\qty{3.0857e18}{cm})^3 = \qty{2.938e55}{cm^3}$: de
verhouding is \qty{37.96}{GeV.cm^{-3}}. Dan $0.30/37.96 = \num{0.0079}$
zonsmassa's per kubieke parsec. (c)~De dichtheid uit de vlakke \omterm{def:b3:milky-way-galaxies:rotation-curve}{rotatiekromme} telt alle massa die
de banen in stand houdt, sterren en gas inbegrepen, alsof die over een
bol verdeeld was; de donkere materie is maar een deel ervan, hier ongeveer de helft.
\end{solution}

\begin{solution}{exo:b3:milky-way-galaxies:7}
(a)~$W_I(\mathrm{LMC}) = -3.314\times1.4771 + 15.888 = \qty{10.993}{mag}$,
dus $M_W = 10.993 - 18.477 = \qty{-7.484}{mag}$. (b)~$\mu = 25.00 + 7.484 =
32.484$ en $d = 10^{32.484/5 + 1}~\mathrm{pc} = \qty{31.4}{Mpc}$.
(c)~$\delta d/d = (\ln 10/5)\,\delta\mu = 0.4605\times0.077 =
\qty{3.5}{\percent}$ voor één ster, $3.5/\sqrt{50} = \qty{0.5}{\percent}$
voor vijftig, als hun fouten onafhankelijk zijn (de fout in de eigen
afstand van de Wolk telt bij alle sterren gelijk op).
\end{solution}

\begin{solution}{exo:b3:milky-way-galaxies:8}
$\mu = I - A_I - M_I = 26.00 - 0.05 + 4.05 = 30.00$, dus $d =
10^{7}~\mathrm{pc} = \qty{10.0}{Mpc}$. Een fout van \qty{0.04}{mag} is
$0.4605\times0.04 = \qty{1.8}{\percent}$ in de afstand.
\end{solution}

\begin{solution}{exo:b3:milky-way-galaxies:9}
(a)~$d = cz/H_0$ levert $5\log_{10}d = 5\log_{10}(cz) - 5\log_{10}H_0$,
dus $m_B - M_B - 25 = 5\log_{10}(cz) - 5\log_{10}H_0$; met
$\log_{10}(cz) = a_B + 0.2\,m_B$, $5\log_{10}H_0 = 5a_B + M_B + 25$, wat
het resultaat is. (b)~$\log_{10}H_0 = -3.8506 + 0.714158 + 5 = 1.8636$, dus
$H_0 = \qty{73.04}{km.s^{-1}.Mpc^{-1}}$. (c)~$M_B = 5(\log_{10}67.4 -
5.714158) = \qty{-19.427}{mag}$: de supernova's zijn \qty{0.174}{mag} helderder,
daarom zijn alle afstanden met een factor $10^{0.2\times0.174} = 1.084$ groter, de
verhouding $73.04/67.4$. De \omterm{def:b3:milky-way-galaxies:tension}{Hubble-spanning} is een verschil van acht procent in de
schaal van afstanden.
\end{solution}

\begin{solution}{exo:b3:milky-way-galaxies:10}
(a)~Een ster aan de rand, in een cirkelvormige baan, heeft $V^2/R \approx GM/R^2$
(sferische benadering). (b)~$R = \sqrt{M/(\pi\Sigma)}$, dus $V^2 =
GM/R = G\sqrt{\pi\Sigma}\,M^{1/2}$ en $M = V^4/(\pi\Sigma G^2) \propto
V^4$; met een gemeenschappelijke massa-lichtverhouding, $L \propto V^4$ ook.
(c)~$(100/229)^4 = \num{0.036}$: het snellere sterrenstelsel is 27 keer
massiever. (d)~Bij gelijke lichtkracht daalt de flux als $1/d^2$, dus de zwakkere
is twee keer zo ver weg.
\end{solution}

\begin{solution}{exo:b3:milky-way-galaxies:11}
(a)~$\qty{1}{km.s^{-1}} = \qty{1.0227}{kpc.Gyr^{-1}}$, dus
\qty{229}{km.s^{-1}} is \qty{234.2}{kpc.Gyr^{-1}}. (b)~$2\pi R/V$:
\qty{107}{Myr} bij \qty{4}{kpc}, \qty{322}{Myr} bij \qty{12}{kpc}.
(c)~$\Delta\phi = Vt\,(1/4 - 1/12) = 234.2\times0.1667 = \qty{39.0}{rad}$,
6.2 omwentelingen: de arm zou een spiraal van zes strakke windingen tussen 4 en
\qty{12}{kpc}, veel strakker dan waargenomen armen, en nog strakker bij elke
verdere rotatie. (d)~In een dichtheidsgolf draait het patroon met één
hoeksnelheid, als een star lichaam; de sterren en het gas passeren het, dus de vorm
van de armen wikkelt zich niet op.
\end{solution}

\begin{solution}{exo:b3:milky-way-galaxies:12}
(a)~$a = (4.2\times10^6\times16^2)^{1/3} = \qty{1024}{au}$; bij
\qty{8178}{pc} spant het $1024/8178 = \ang{;;0.125}$ (een au op één parsec
spant één boogseconde). (b)~$R_S = 2GM/c^2 = 2\times\num{6.674e-11}
\times\num{8.35e36}/(\num{2.998e8})^2 = \qty{1.24e10}{m} = \qty{0.083}{au}$;
de diameter $2R_S$ spant $0.166/8178 = \num{2.0e-5}$ boogseconde, ongeveer
\qty{20}{\micro as}. (c)~$R_S = \qty{4.02e-7}{pc}$, en $M/(\tfrac43\pi
R_S^3) = \num{1.5e25}$ zonsmassa's per kubieke parsec, veertig keer de
ondergrens die door de beweging en de grootte van Sgr~A* wordt vastgelegd: de grens komt binnen
een factor veertig van de dichtheid van het zwarte gat zelf.
\end{solution}

\begin{solution}{pb:b3:milky-way-galaxies:1}

\textbf{Deel I.}

\textbf{1.} $v_r = A\,d\sin2\ell$: $+A d$ bij \ang{45}, 0 bij \ang{90}, $-A d$ bij
\ang{135}, 0 bij \ang{180}.

\textbf{2.} $A = (15.4 + 15.2)/(2\times\qty{1.00}{kpc}) =
\qty{15.3}{km.s^{-1}.kpc^{-1}}$.

\textbf{3.} Bij \ang{90}, $v_t = d(B - A) = 4.74\times(-5.91)\times1.00 =
\qty{-28.02}{km.s^{-1}}$, dus $B = 15.3 - 28.02 =
\qty{-12.72}{km.s^{-1}.kpc^{-1}}$.

\textbf{4.} Bij \ang{0}, $v_t = d(A + B) = 4.74\times0.55 =
\qty{2.61}{km.s^{-1}}$, dus $B = 2.61 - 15.3 = \num{-12.69}$: de twee
stemmen overeen binnen \num{0.03}.

\textbf{5.} $\omega_0 = A - B = \qty{28.0}{km.s^{-1}.kpc^{-1}}$ en
$V_0 = 28.0\times8.18 = \qty{229}{km.s^{-1}}$.

\textbf{6.} $V'_0 = -(A + B) = \qty{-2.6}{km.s^{-1}.kpc^{-1}}$: langzaam dalend.

\textbf{7.} $P = 2\pi/\omega_0$: $2\pi\times\qty{8.18}{kpc}/\qty{229.2}{km.s^{-1}}
= \qty{219}{Myr}$; $4540/219 = 20.7$ omlopen.

\textbf{Deel II.}

\textbf{8.} Het raakpunt vereist een zichtlijn die dichter bij het
centrum passeert dan waar de Zon staat, $\sin\ell > 0$ met $\ell < \ang{90}$ (of het
spiegelkwadrant, $\ang{270} < \ell < \ang{360}$, waar de grootste
snelheid het meest negatief is). Buiten de zonnecirkel neemt de
afstand tot het centrum langs de hele zichtlijn toe: geen punt springt
eruit, en de snelheid alleen zegt niet waar het gas zich bevindt.

\textbf{9.} $R = 8.18\sin\ell$: \qtylist{5.26;6.27;7.08;7.69}{kpc}.

\textbf{10.} $V = v_\mathrm{max} + V_0\sin\ell$: $82.9 + 147.2 = 230.1$,
$55.6 + 175.4 = 231.0$, $32.4 + 198.3 = 230.7$, $14.9 + 215.2 =
\qty{230.1}{km.s^{-1}}$.

\textbf{11.} $\Delta\nu = \nu v/c = \num{1.420406e9}\times82.9/299792 =
\qty{393}{kHz}$, naar lagere frequenties: het gas beweegt zich van ons af.

\textbf{12.} Binnen ongeveer \qty{3}{km.s^{-1}} van de rode reuzen: tussen 5 en
\qty{8}{kpc} is de kromme vlak.

\textbf{13.} Een puntmassa geeft $V = V_0\sqrt{R_0/R} = 229\times\sqrt{8.18/5.26} =
\qty{286}{km.s^{-1}}$ bij het eerste raakpunt, veel sneller dan de
\qty{230}{km.s^{-1}} gevonden: de massa is niet in het centrum geconcentreerd,
een groot deel ervan ligt tussen 5 en \qty{8}{kpc} en verder.

\textbf{Deel III.}

\textbf{14.} $\num{1.067e10} + 2\times\num{3.944e10} = \num{8.96e10}$ zonsmassa's,
bijna twee keer de \num[separate-uncertainty=true]{5 \pm 1 e10} uit de review; de sterren van het model
zijn mogelijk te zwaar, waardoor de hieronder gevonden donkere materie
een onderschatting zou zijn.

\textbf{15.} $V_b^2 = GMR^2/(R^2 + b^2)^{3/2} = \num{4.301e-6}\times\num{1.067e10}
\times400/(400.09)^{3/2}$, $V_b = \qty{47.9}{km.s^{-1}}$.

\textbf{16.} Dunne schijf: $(a + b)^2 = 30.80$, $V^2 = \num{4.301e-6}\times\num{3.944e10}
\times400/(430.8)^{3/2}$, $V = \qty{87.1}{km.s^{-1}}$. \omterm{def:b3:milky-way-galaxies:structure}{Dikke schijf}:
$(a + b)^2 = 11.56$, $V = \qty{90.1}{km.s^{-1}}$.

\textbf{17.} $V_\mathrm{vis} = \sqrt{47.9^2 + 87.1^2 + 90.1^2} =
\qty{134.2}{km.s^{-1}}$.

\textbf{18.} Bij \qty{8.18}{kpc}: 74.8, 108.4 en \qty{127.8}{km.s^{-1}},
samen \qty{183.5}{km.s^{-1}}, tegenover $V_0 = \qty{229}{km.s^{-1}}$:
de zichtbare massa schiet al bij de Zon tekort.

\textbf{19.} $\sqrt{G\times\num{8.96e10}/20} = \qty{138.8}{km.s^{-1}}$. Het
model komt met \qty{134.2}{km.s^{-1}} dicht in de buurt, omdat op \qty{20}{kpc},
vier tot zes schijfschaallengten verderop, ligt vrijwel alle zichtbare massa
binnen de baan: de zichtbare kromme is keplerisch geworden.

\textbf{Deel IV.}

\textbf{20.} $M = 20\times200^2/\num{4.301e-6} = \num{1.860e11}$ zonsmassa's.

\textbf{21.} $M_\mathrm{vis} = 20\times134.2^2/\num{4.301e-6} = \num{8.37e10}$
zonsmassa's.

\textbf{22.} $M_\mathrm{dark} = \num{1.860e11} - \num{8.37e10} = \num{1.02e11}$;
aangezien beide massa's $RV^2/G$ bij dezelfde $R$ zijn, is de fractie
$1 - (V_\mathrm{vis}/V)^2 = 1 - (134.2/200)^2 = \num{0.550}$.

\textbf{23.} $x = 20/14.8 = 1.351$, $\ln(2.351) - 1.351/2.351 = 0.8549 - 0.5748
= 0.2801$, $M = 4\pi\times\num{1.06e7}\times14.8^3\times0.2801 =
\num{1.21e11}$ zonsmassa's; de donkere fractie is
$1.21/(1.21 + 0.837) = \num{0.59}$. De halo is gefit op de hele
kromme, die bij \qty{20}{kpc} doorschiet: de metingen zakken daar
met ongeveer \qty{10}{km.s^{-1}} onder het gladde model (\qty{209}{km.s^{-1}}).

\textbf{24.} $x = 8.18/14.8 = 0.553$, $\rho = \num{1.06e7}/(0.553\times1.553^2)
= \num{7.95e6}$ zonsmassa's per kubieke kiloparsec, \num{0.0080} per
kubieke parsec, of \qty{0.30}{GeV.cm^{-3}}, de waarde die de auteurs van de
fit opgeven.

\textbf{25.} $1 - 8.37/18.60 = \num{0.550}$: de fractie donkere materie van de
massa binnen \qty{20}{kpc} is
\textbf{\qty[text-series-to-math=true]{55}{\percent}}.
\end{solution}
